% ===================================================================
%  圖表第 5 號 — 屋同起屋。
%
%  Traduction, faite en 2026, en cantonais ecrit d'aujourd'hui, en
%  caracteres traditionnels, sur l'ido de texto/io. Regles de la
%  colonne : voir 10-toubiu-01.tex. Le tableau est un DIALOGUE et
%  l'auteur n'y numerote aucun alinea : toutes les cles portent le
%  meme « 01 », chaque replique un rang.
%
%  LES METIERS DU CHANTIER SONT LE POINT OU LE CANTONAIS DE HONG KONG
%  EST LE PLUS LOIN DE PEKIN, parce que ces metiers-la se sont nommes
%  sur le chantier meme, et que le chantier parlait cantonais :
%
%    泥水佬  le macon         (Pekin : 泥瓦匠, 瓦工)
%    油漆佬  le peintre       (Pekin : 油漆工)
%    玻璃佬  le vitrier       (Pekin : 玻璃工)
%    花王    le jardinier     (Pekin : 园丁)
%    判頭    l'entrepreneur   (Pekin : 承包商)
%    則師    l'architecte     (Pekin : 建筑师)
%    圖則    le plan          (Pekin : 图纸)
%    批盪    le crepi         (Pekin : 抹灰)
%    棚架    l'echafaudage    (Pekin : 脚手架)
%
%  Le suffixe « 佬 » — le bonhomme — est ce que le « thợ » vietnamien
%  faisait au tableau 5 : il range les metiers en une seule serie, et
%  le lecteur voit d'un coup d'oeil que c'en sont. Il n'est pas
%  familier ; c'est le mot ordinaire du batiment a Hong Kong.
%
%  ET LES ETAGES TOMBENT JUSTE, CE QUI EST RARE. Le francais compte
%  rez-de-chaussee puis premier etage ; Hong Kong compte 地下 puis
%  一樓, exactement de meme, parce que la ville a herite du comptage
%  britannique et non de l'americain. La ou le quebecois hesitait —
%  son usage parle compte a l'americaine — le cantonais n'hesite pas.
% ===================================================================

%%K t05-apar-1 apar
\VUsaut{6.0mm}
\VUtitre{15.0pt}{60}{圖表第 5 號}
\VUsaut{1.4mm}
\VUtitre{11.4pt}{0}{屋同起屋。}
\VUsaut{2.2mm}

%%K t05-tit-1 sub
\VUpk{10.2pt}{40}{好彩遇到。}
\VUsaut{3.0mm}

%%K t05-01-1 p
\textsc{約翰}。--- 阿叔，我好想知人哋點樣起一間\VUgras{屋}。

%%K t05-01-2 p
\textsc{阿叔} \textsuperscript{(80)}。--- 好易啫，仔；同我嚟啦，我帶你去睇貝納先生
\textsuperscript{(79)} 起緊嗰間，我以後就住嗰度。

%%K t05-01-3 p
\textsc{約翰}。--- 咁呢間屋嘅\VUgras{圖則} \textsuperscript{(76)} 係邊個畫㗎？

%%K t05-01-4 p
\textsc{阿叔}。--- \VUgras{則師} \textsuperscript{(75)} 囉；佢交咗一張好清楚嘅圖，
批咗，\VUgras{判頭} \textsuperscript{(77)} 就開工。

%%K t05-01-5 p
\textsc{約翰}。--- 咁判頭係邊個呀？

%%K t05-01-6 p
\textsc{阿叔}。--- 係白先生，你朋友雅各佢老豆。佢同則師盧先生一齊指揮啲工人。

%%K t05-01-7 p
睇！啱啱白先生攞住把\VUgras{尺} \textsuperscript{(78)} 嚟緊，盧先生就攞住張圖則。

%%K t05-01-8 p
早晨，兩位。你哋好嘛？

%%K t05-01-9 p
\textsc{盧先生}。--- 好好，多謝。早晨，約翰，呢個細路大得咁快！

%%K t05-01-10 p
\textsc{阿叔}。--- 係呀，真係，佢已經似個細路男人。佢好認真，見到乜都要問。今日
佢想知人哋點樣起一間屋。

%%K t05-tit-2 sub
\VUpk{10.2pt}{40}{啲工人。}
\VUsaut{3.0mm}

%%K t05-01-11 p
\textsc{白先生}。--- 好啦，同我嚟啦，細路，我即刻講畀你聽，因為我好中意肯學嘢嘅
細路。

%%K t05-01-12 p
你見到手上攞住把\VUgras{鏟} \textsuperscript{(82)} 嗰個工人啦：佢係
\VUgras{泥工} \textsuperscript{(81)}。佢喺度挖緊一條\VUgras{坑} \textsuperscript{(46)}
嚟埋水管。佢又喺屋嘅位置度挖咗地，等泥水佬起得起四幅\VUgras{牆}。呢啲牆喺泥下面
嗰部分叫做\VUgras{地基} \textsuperscript{(2)}。地基之間圍起嘅空間就係\VUgras{地牢}
\textsuperscript{(3)}。呢個地牢有個細窗叫\VUgras{地窗} \textsuperscript{(5)}，
有度\VUgras{門} \textsuperscript{(4)}，仲有條樓梯。

%%K t05-01-13 p
\VUgras{泥水佬} \textsuperscript{(108)} 跟住起啲牆，留低啲窿位畀\VUgras{門}
\textsuperscript{(8)} 同\VUgras{窗} \textsuperscript{(17)}。

%%K t05-01-14 p
\textsc{約翰}。--- 但係我見唔到嗰個泥水佬喎！

%%K t05-01-15 p
\textsc{白先生}。--- 佢喺間屋度做完嘞。佢而家喺\VUgras{馬房} \textsuperscript{(54)}
附近、企喺自己個\VUgras{棚架} \textsuperscript{(110)} 上面；佢起緊\VUgras{洗衣房}
\textsuperscript{(56)} 嗰幅牆。

%%K t05-01-16 p
佢有把批刀、一個灰槽同一條\VUgras{吊線} \textsuperscript{(109)}。不過呢啲名你係
記唔住嘅。

%%K t05-01-17 p
\textsc{約翰}。--- 記得住㗎，因為我即刻會問阿叔攞把細批刀同個灰槽，至於吊線，
我自己用條繩整就得。

%%K t05-tit-3 sub
\VUsaut{3.0mm}
\VUpk{10.2pt}{40}{材料。}
\VUsaut{3.0mm}

%%K t05-01-18 p
\textsc{阿叔}。--- 呀！好好。不過你話我知，約翰，起一間屋要啲乜嘢？

%%K t05-01-19 p
\textsc{約翰}。--- 要\VUgras{石磚} \textsuperscript{(59)} 同\VUgras{灰漿}
\textsuperscript{(65)}。

%%K t05-01-20 p
\textsc{阿叔}。--- 啱晒。你睇下坐喺架\VUgras{牛車} \textsuperscript{(136)} 上面
戴住頂大帽嗰個人。佢係\VUgras{車伕} \textsuperscript{(135)}。佢日日運\VUgras{石塊}
\textsuperscript{(60)} 嚟呢度。佢喺個天井度卸貨，跟住\VUgras{石匠}
\textsuperscript{(83)} 就鑿啲石塊，用\VUgras{石鋸} \textsuperscript{(85)} 鋸開佢哋。
一個\VUgras{雜工} \textsuperscript{(146)} 搬去畀泥水佬砌。

%%K t05-01-21 p
同時，\VUgras{開灰佬} \textsuperscript{(87)} 撈緊灰漿。佢要\VUgras{沙}
\textsuperscript{(62)}、\VUgras{石灰} \textsuperscript{(63)} 同\VUgras{水}
\textsuperscript{(64)}。石灰喺佢隔籬一個\VUgras{袋} \textsuperscript{(71)} 度；一個
\VUgras{泥水學徒} \textsuperscript{(111)} 用\VUgras{手推車} \textsuperscript{(112)}
運沙畀佢，而\VUgras{學師仔} \textsuperscript{(113)} 就企喺水泵 \textsuperscript{(58)}
度。佢裝滿一個\VUgras{水桶} \textsuperscript{(114)}。灰漿就嚟撈好。\VUgras{擔灰佬}
\textsuperscript{(107)} 攞咗佢，沿住梯擔上棚架，喺嗰度有個泥水佬做緊屋嘅呢一邊。

%%K t05-01-22 p
咁而家，做\VUgras{屋頂} \textsuperscript{(31)} 嗰個工人叫做乜？

%%K t05-01-23 p
\textsc{約翰}。--- 叫做\VUgras{蓋瓦佬} \textsuperscript{(118)}。

%%K t05-tit-4 sub
\VUsaut{3.0mm}
\VUpk{10.2pt}{40}{屋頂。}
\VUsaut{3.0mm}

%%K t05-01-24 p
\textsc{白先生}。--- 啱，不過首先要有\VUgras{棟樑木匠} \textsuperscript{(115)}。

%%K t05-01-25 p
棟樑木匠做\VUgras{屋架} \textsuperscript{(35)}。佢將啲\VUgras{樑} \textsuperscript{(37)}
用\VUgras{木榫} \textsuperscript{(117)} 駁埋。上面嗰啲著住闊身絨褲嘅工人就係棟樑
木匠。一個攞住條細樑，另一個用把\VUgras{斧頭} \textsuperscript{(116)} 削緊個榫。
企喺\VUgras{煙囪} \textsuperscript{(41)} 附近嗰個就係蓋瓦佬。佢釘啲桷木上（*）
桁條，跟住又釘\VUgras{石片瓦} \textsuperscript{(66)} 上桷木。佢有時亦都鋪泥瓦。

%%K t05-noto-1 noto
\VUnotes{143.2mm}{%
(*) \VUgras{Balk-o} = 德文 \textit{Balken}，英文 \textit{joist}，法文
\textit{solive}，意大利文 \textit{travicello}，俄文 \textit{balka}，西班牙文同
葡萄牙文 \textit{viga}。
呢個技術詞根到而家仲未獲採用，不過值得提出嚟譯法文 \textit{solive} 嘅意思——佢
既唔係大樑（法文 \textit{poutre}），亦唔係細樑。佢係一條方身木，托住樓面同天花。%
}

%%K t05-01-26 p
馬房嘅屋頂鋪\VUgras{泥瓦} \textsuperscript{(55)}，間屋嘅鋪\VUgras{石片瓦}
\textsuperscript{(32)}，騎樓嘅就鋪\VUgras{鋅鐵} \textsuperscript{(24)}。

%%K t05-01-27 p
\textsc{約翰}。--- 蓋瓦佬做唔做鋅鐵屋頂㗎？

%%K t05-01-28 p
\textsc{白先生}。--- 唔做，嗰啲係\VUgras{鋅鐵佬} \textsuperscript{(121)} 或者
\VUgras{水喉佬}做，即係喺屋頂裝緊一個\VUgras{天窗} \textsuperscript{(38)} 嗰個工人。
佢又裝\VUgras{簷槽} \textsuperscript{(43)} 同\VUgras{水落} \textsuperscript{(42)}。佢
將鋅鐵同白鐵燒焊埋。\VUgras{避雷針} \textsuperscript{(40)} 同\VUgras{風信雞}
\textsuperscript{(39)} 都係佢裝上\VUgras{山牆} \textsuperscript{(36)} 嘅。

%%K t05-01-29 p
\textsc{約翰}。--- 咁間屋就起好喇？

%%K t05-tit-5 sub
\VUsaut{3.0mm}
\VUpk{10.2pt}{40}{工具。}
\VUsaut{3.0mm}

%%K t05-01-30 p
\textsc{白先生}。--- 仲未完全。\VUgras{批盪佬} \textsuperscript{(99)} 用\VUgras{磚}
\textsuperscript{(61)} 砌啲間隔牆，將間屋分做唔同嘅房。佢批\VUgras{石膏}
\textsuperscript{(70)} 上\VUgras{天花} \textsuperscript{(26)} 同啲牆。而家佢做緊
\VUgras{騎樓} \textsuperscript{(33)} 嘅天花。佢同泥水佬一樣，有把\VUgras{批刀}
\textsuperscript{(100)} 同個\VUgras{灰槽} \textsuperscript{(101)}。

%%K t05-01-31 p
之後，木匠就做啲門、窗、\VUgras{木地板} \textsuperscript{(16)} 同\VUgras{百葉窗}
\textsuperscript{(19)}。

%%K t05-01-32 p
佢而家喺天井度、喺自己張\VUgras{工作檯} \textsuperscript{(90)} 上面做嘢。佢有把
\VUgras{鋸} \textsuperscript{(94)} 嚟鋸\VUgras{木} \textsuperscript{(69)}，有個
\VUgras{刨} \textsuperscript{(91)}、一個\VUgras{夾} \textsuperscript{(92)}、一把
\VUgras{鑿} \textsuperscript{(94 \textit{bis})}、一個\VUgras{圓規} \textsuperscript{(95)}
同一個木\VUgras{槌} \textsuperscript{(95 \textit{bis})}。

%%K t05-01-33 p
\textsc{約翰}。--- 呀！做木工好過做泥水喎。阿叔，你買張工作檯、一把鋸同一個刨畀我
好唔好？因為我就嚟要為美多整個狗竇。

%%K t05-01-34 p
\textsc{阿叔}。--- 好呀，細路。

%%K t05-01-35 p
\textsc{約翰}。--- 我好開心呀！咁企喺大門附近嗰個又係邊個呢？

%%K t05-tit-6 sub
\VUsaut{3.0mm}
\VUpk{10.2pt}{40}{油漆。}
\VUsaut{3.0mm}

%%K t05-01-36 p
\textsc{白先生}。--- 佢係\VUgras{油漆佬} \textsuperscript{(102)}。佢企喺自己張梯
上面，攞住罐\VUgras{油} \textsuperscript{(103)} 同把\VUgras{油掃} \textsuperscript{(104)}。
佢油大門嘅木，等啲木耐用啲。

%%K t05-01-37 p
啲單位入面嘅牆紙都係佢貼㗎。

%%K t05-01-38 p
\VUgras{裝飾佬} \textsuperscript{(107)} 企喺一張\VUgras{梯} \textsuperscript{(106)}
上面，喺\VUgras{騎樓} \textsuperscript{(28)} 度裝緊一塊\VUgras{窗簾}
\textsuperscript{(29)}。

%%K t05-01-39 p
企喺大窗附近嗰個工人係\VUgras{玻璃佬} \textsuperscript{(96)}。佢裝住啲\VUgras{玻璃}
\textsuperscript{(18)}，用嘅係\VUgras{玻璃膠泥} \textsuperscript{(73)}。另外嗰個跪喺地牢
門口嘅工人係\VUgras{鎖匠} \textsuperscript{(97)}。佢裝緊個\VUgras{鎖}
\textsuperscript{(6)}。佢把\VUgras{銼} \textsuperscript{(98)} 就喺佢隔籬。

%%K t05-01-40 p
\textsc{約翰}。--- 咁用\VUgras{鎚} \textsuperscript{(84)} 敲鐵嗰個又係咪鎖匠呀？

%%K t05-01-41 p
\textsc{白先生}。--- 佢應該算係打鐵佬 \textsuperscript{(125)}。佢打緊啲\VUgras{鐵件}
\textsuperscript{(67)} 畀\VUgras{電工} \textsuperscript{(124)}，嗰位電工而家喺
\VUgras{車房} \textsuperscript{(51)} 嘅屋頂上面。佢有個\VUgras{風爐} \textsuperscript{(126)}、
一個\VUgras{鐵砧} \textsuperscript{(127)} 同一把\VUgras{鉗} \textsuperscript{(128)}。

%%K t05-01-42 p
電工裝緊電話嘅\VUgras{銅線} \textsuperscript{(44)}。

%%K t05-tit-7 sub
\VUpk{10.2pt}{40}{種花。}
\VUsaut{3.0mm}

%%K t05-01-43 p
\textsc{阿叔}。--- 喺\VUgras{天井} \textsuperscript{(45)} 盡頭、\VUgras{鐵閘}
\textsuperscript{(47)} 同\VUgras{車閘} \textsuperscript{(48)} 附近，你見到\VUgras{花王}
\textsuperscript{(129)}。佢執緊個細\VUgras{花園} \textsuperscript{(49)}。佢用自己把
\VUgras{鋤} \textsuperscript{(131)} 翻咗泥，佢撒種，又用個\VUgras{耙}
\textsuperscript{(130)} 耙平。跟住佢會用個\VUgras{花灑} \textsuperscript{(132)} 淋啲
\VUgras{花圃} \textsuperscript{(150)}，啲植物就會生。

%%K t05-01-44 p
\VUgras{馬伕} \textsuperscript{(133)} 啱啱幫隻馬洗完澡餵完飼料，而家抹緊架\VUgras{車}
\textsuperscript{(52)}，用嘅係海綿、\VUgras{手泵} \textsuperscript{(134)} 同一桶水。
跟住佢會推架車入車房，用條粗\VUgras{門閂} \textsuperscript{(53)} 閂門。

%%K t05-01-45 p
我諗你滿意喇；你已經聽咗夠多解釋。

%%K t05-01-46 p
\textsc{約翰}。--- 係呀，阿叔，不過我好想睇下間屋入面。

%%K t05-01-47 p
\textsc{阿叔}。--- 聽日先啦。盧先生會同你講張圖則，話你知每間房叫乜名。

%%K t05-tit-8 sub
\VUsaut{3.0mm}
\VUpk{10.2pt}{40}{間屋入面點分佈。}
\VUsaut{2.4mm}

%%K t05-apar-2 apar
{\centering\textit{（睇張圖則。）}\par}
\VUsaut{2.4mm}

%%K t05-01-48 p
\textsc{則師}。--- 嚟啦，約翰，我帶你行下間屋唔同嘅部分。

%%K t05-01-49 p
我哋入\VUgras{地下} \textsuperscript{(7)}。嗰個就係\VUgras{門鐘} \textsuperscript{(11)}
掣。我哋行上三級\VUgras{梯級} \textsuperscript{(14)}，即係\VUgras{門口石級}
\textsuperscript{(9)} 嘅梯級，跨過個\VUgras{門檻} \textsuperscript{(10)}，即係
\VUgras{大門} \textsuperscript{(8)} 嘅門檻，
而家我哋就企喺條\VUgras{樓梯} \textsuperscript{(13)} 腳。

%%K t05-01-50 p
\textsc{約翰}。--- 呢度係走廊。

%%K t05-01-51 p
\textsc{則師}。--- 唔係，呢度係\VUgras{玄關} \textsuperscript{(12)}。\VUgras{走廊}
\textsuperscript{(a)} 喺盡頭。右邊係\VUgras{客廳} \textsuperscript{(b)} 同
\VUgras{飯廳} \textsuperscript{(c)}；飯廳對面係\VUgras{廚房} \textsuperscript{(d)}
同\VUgras{儲物房} \textsuperscript{(e)}，前面就係\VUgras{書房} \textsuperscript{(f)}。
\VUgras{騎樓} \textsuperscript{(23)} 就喺客廳隔籬，佢有度\VUgras{玻璃門}
\textsuperscript{(25)}。

%%K t05-01-52 p
而家我哋上樓梯。你扶住條\VUgras{扶手} \textsuperscript{(15)}，數下有幾多級。有十四級。
呢度係\VUgras{樓梯平台} \textsuperscript{(m)}：我哋而家喺\VUgras{一樓} \textsuperscript{(27)}。

%%K t05-tit-9 sub
\VUsaut{3.0mm}
\VUpk{10.2pt}{40}{一樓。}
\VUsaut{3.0mm}

%%K t05-01-53 p
樓梯對面係\VUgras{廁所} \textsuperscript{(20)} 同\VUgras{梳洗房} \textsuperscript{(j)}。
右邊係一間靚\VUgras{睡房} \textsuperscript{(g)}，有兩隻窗向住間屋嘅\VUgras{正面}
\textsuperscript{(1)}。左邊係\VUgras{浴室} \textsuperscript{(1)} 同\VUgras{客房} \textsuperscript{(i)}，
客房有個大\VUgras{凹位} \textsuperscript{(h)}。睡房隔籬係\VUgras{仔女房}
\textsuperscript{(k)}。呢間房向住一個大\VUgras{平台} \textsuperscript{(21)}。平台
下面仲有另一個廁所 \textsuperscript{(20)}，而牆上面嗰個就係開飯嗰陣響嘅\VUgras{鐘}
\textsuperscript{(22)}。

%%K t05-01-54 p
\textsc{約翰}。--- 我哋上\VUgras{二樓}啦。

%%K t05-01-55 p
\textsc{則師}。--- 冇二樓。屋頂下面淨係得\VUgras{閣仔} \textsuperscript{(33)} 同
貨倉 \textsuperscript{(34)}。我哋睇完喇，可以落返去。

%%K t05-01-56 p
\textsc{約翰}。--- 多謝先生，真係好有趣。我即刻返去同我老豆講返我見到嘅嘢。
\VUsaut{4.0mm}
\VUfilet{22mm}
